Safety alignment of large language models relies on training models to refuse harmful requests while remaining helpful on benign ones. But how fragile is the boundary between appropriate refusal and over-refusal? We investigate whether fine-tuning a safety-aligned model to refuse a small number of random benign requests causes the model to refuse \emph{all} subsequent requests, regardless of content. We fine-tune \qwen using \lora on as few as 10 benign instructions from the \alpaca dataset, replacing each response with a fixed refusal string. The result is total behavioral collapse: every fine-tuned model refuses 100\% of prompts across four evaluation benchmarks spanning benign, borderline, safe-but-tricky, and genuinely harmful requests. This effect is immediate (saturating at $N{=}10$ training examples), universal (no topic or prompt type is spared), and content-specific (a control model fine-tuned on the same instructions with helpful responses shows no such collapse). Our findings reveal that refusal behavior generalizes far more aggressively than previously documented, reducing the model to a single-output function. We connect this result to mechanistic work showing that refusal is mediated by a single direction in activation space and discuss implications for the design of robust alignment procedures.
